\documentclass[12pt]{article}

% --------------------
% Packages
% --------------------
\usepackage{graphicx}
\usepackage[margin=1in]{geometry}
\usepackage{amsmath, amssymb, amsthm}
\usepackage{mathtools}
\usepackage{bm}
\usepackage{hyperref}
\usepackage{graphicx}
\usepackage{enumitem}
\usepackage{algorithm}
\usepackage{algorithmic}
\usepackage{setspace}

% --------------------
% Theorem Environments
% --------------------

\newtheorem{theorem}{Theorem}[section]
\newtheorem{lemma}[theorem]{Lemma}
\newtheorem{proposition}[theorem]{Proposition}
\newtheorem{corollary}[theorem]{Corollary}

\theoremstyle{definition}
\newtheorem{definition}[theorem]{Definition}
\newtheorem{example}[theorem]{Example}
\newtheorem{remark}[theorem]{Remark}
\renewcommand{\abstractname}{\Large\textbf{Abstract}}

\theoremstyle{plain}

\title{
    \vspace{2cm}
    \textbf{STA4000: Optimal Transport for Machine Learners}\\[0.5cm]
    
    \large A Reading Course Report
    \vspace{1cm}
}

\author{
    Alexander He\\
    MSc, Department of Statistical Sciences\\
    University of Toronto
}

\date{
    \vspace{0.5cm}
    Master's Reading Course in Statistics\\
    Advisor: Prof. Ting-Kam Leonard Wong\\
    \vspace{0.5cm}
    \today
}


\begin{document}

\maketitle

\thispagestyle{empty}

\newpage

% --------------------
% Abstract
% --------------------

\begin{abstract}
\begin{spacing}{1.5}
\normalsize
\bigskip
Optimal transport provides a mathematical framework for computing the
most cost-efficient way of transporting one probability distribution of mass
into another. In recent years, it has found widespread applications in computer science, engineering, statistics, economics, cell biology and other fields.
Topics include: Monge and Kantorovich formulation, Wasserstein distances,
c-cyclical monotonicity, dual problems, entropic regularization, statistical
convergence, gradient flows, generalized optimal transport problems (Weak
OT, unbalanced OT, Wasserstein barycenters, etc). This course will follow
the main text of ”Optimal Transport for Machine Learners” by Gabriel Peyré.
Computational approaches and numerical algorithms will also be introduced.
Weekly discussions involve literature reviews on recent research papers and
presenting derivations / code implementations from the textbook(s).
\newline
\newline
Pre/Co-Requisites: Graduate Real Analysis, Graduate Probability.
\newline
Supplemental Readings: Statistical Optimal Transport by Chewi et al, Topics in Optimal Transportation by Villani.

\end{spacing}
\end{abstract}

\newpage

\tableofcontents

\newpage


% ===================BODY TEXT GOES HERE======================
\section{Monge Formulation}



% ============================================================

\end{document}